% ============================================================
%  E-Mualim CRAG Conference Paper — IEEE Format
%  Compile: pdflatex -> bibtex -> pdflatex -> pdflatex
%  Template: IEEEtran (standard IEEE conference)
% ============================================================

\documentclass[conference]{IEEEtran}

% ---------- Core packages ----------
\usepackage{cite}
\usepackage{amsmath,amssymb,amsfonts}
\usepackage{graphicx}
\usepackage{textcomp}
\usepackage{xcolor}
\usepackage{booktabs}
\usepackage{multirow}
\usepackage{url}
\usepackage{microtype}
\usepackage{subcaption}
\usepackage{array}
\usepackage{tabularx}

% ---------- TODO helper (remove before final submission) ----------
\newcommand{\todo}[1]{\textcolor{red}{\textbf{[TODO: #1]}}}

% ============================================================
\begin{document}

\title{Corrective Retrieval-Augmented Generation for\\
Hallucination Mitigation in Educational Q\&A:\\
A Case Study on E-Mualim}

% IEEE conference — list all authors
\author{
  \IEEEauthorblockN{Author One}
  \IEEEauthorblockA{
    \textit{Department of Computer Science}\\
    \textit{FAST-NUCES Lahore}\\
    Lahore, Pakistan\\
    authorone@nu.edu.pk
  }
  \and
  \IEEEauthorblockN{Author Two}
  \IEEEauthorblockA{
    \textit{Department of Computer Science}\\
    \textit{FAST-NUCES Lahore}\\
    Lahore, Pakistan\\
    authortwo@nu.edu.pk
  }
  \and
  \IEEEauthorblockN{Author Three}
  \IEEEauthorblockA{
    \textit{Department of Computer Science}\\
    \textit{FAST-NUCES Lahore}\\
    Lahore, Pakistan\\
    authorthree@nu.edu.pk
  }
}

\maketitle

% ============================================================
%  ABSTRACT
% ============================================================
\begin{abstract}
Large language models (LLMs) deployed in educational tutoring systems
pose a critical safety risk: they generate factually incorrect content
with high confidence, misleading students who lack the background
knowledge to identify such errors.
Standard retrieval-augmented generation (RAG) partially addresses this
by grounding responses in a knowledge base, yet naive RAG passes all
retrieved chunks directly to the generator regardless of their
relevance to the specific query, leaving hallucination risk unmitigated.
We present \textbf{E-Mualim}, a speech-enabled AI tutoring system
that implements a \emph{two-stage corrective pipeline}: a CRAG-style
chunk-level relevance grader that filters retrieved content before
generation, and a post-generation hallucination checker that retries
generation when output is not grounded in source material.
We evaluate E-Mualim against a naive RAG baseline on a curated
educational Q\&A test set using the RAGAS evaluation framework,
reporting faithfulness, answer relevancy, context precision, and
context recall.
\todo{Insert 2--3 key result numbers once evaluation is complete.}
Our results demonstrate that the two-stage corrective architecture
substantially reduces hallucination while preserving answer quality,
providing a practical and reproducible design for safe LLM deployment
in educational settings.
\end{abstract}

\begin{IEEEkeywords}
large language models, hallucination mitigation, retrieval-augmented
generation, corrective RAG, educational AI, intelligent tutoring systems
\end{IEEEkeywords}

% ============================================================
%  I. INTRODUCTION
% ============================================================
\section{Introduction}
\label{sec:intro}

The integration of large language models (LLMs) into educational
technology has created new possibilities for personalised, on-demand
student tutoring.
Systems capable of answering curriculum questions in natural language
can reduce dependency on instructor availability and scale educational
support to large student populations.
However, this promise is fundamentally undermined by a well-documented
failure mode: LLMs \emph{hallucinate}, generating plausible-sounding
but factually incorrect statements with high confidence
\cite{ji2023survey,huang2025survey}.

In general-purpose applications, a hallucinated fact is an
inconvenience.
In an educational tutoring system, it risks permanently embedding a
misconception in a student who has no independent means to verify
the response.
Empirical research has confirmed this concern: erroneous LLM feedback
in tutoring systems has been shown to increase student confusion and
task completion time even when overall learning gains are preserved
\cite{steinbach2025hallucinate}.

Retrieval-Augmented Generation (RAG) \cite{lewis2020rag} has become
the standard approach to grounding LLM responses in verifiable
knowledge sources.
By retrieving relevant documents at query time and conditioning
generation on them, RAG reduces reliance on parametric memory.
However, \emph{naive} RAG offers no quality assurance on the retrieved
context: it passes all top-$k$ retrieved chunks directly to the
generator, including those that are weakly or irrelevantly matched to
the specific query.
Prior work has demonstrated that LLMs are easily distracted by
irrelevant context in the prompt \cite{shi2023distracted}, meaning
that naive RAG can \emph{introduce} hallucination by contaminating
the generation context with off-topic material.

Corrective Retrieval-Augmented Generation (CRAG)
\cite{yan2024crag} addresses this gap by inserting a lightweight
retrieval evaluator between retrieval and generation to filter
low-quality chunks before they reach the generator.
We adapt CRAG for the specific requirements of an educational Q\&A
domain and extend it with a second corrective stage: a
post-generation hallucination checker that detects unfaithful
generation and triggers a retry before the response is delivered to
the student.

The contributions of this paper are as follows:
\begin{itemize}
  \item We present \textbf{E-Mualim}, an educational tutoring system
        implementing a two-stage corrective RAG pipeline that combines
        pre-generation chunk filtering with post-generation faithfulness
        verification.
  \item We describe a practical implementation of the CRAG grader
        using a small LLM (\texttt{llama-3.1-8b-instant}) as a
        batched prompt-based relevance scorer, requiring no
        fine-tuning on labelled retrieval data.
  \item We provide an empirical evaluation of E-Mualim against a
        naive RAG baseline using the RAGAS framework, demonstrating
        measurable hallucination reduction across four complementary
        metrics.
\end{itemize}

% ============================================================
%  II. RELATED WORK
% ============================================================
\section{Related Work}
\label{sec:related}

\subsection{Hallucination in Large Language Models}

Hallucination in natural language generation was formalised by
Ji et al. \cite{ji2023survey}, who distinguish \emph{intrinsic}
hallucination (output that contradicts the source material) from
\emph{extrinsic} hallucination (output containing information absent
from and unverifiable against the source).
Huang et al. \cite{huang2025survey} refine this taxonomy for
LLMs, separating \emph{factuality hallucination} (contradicting world
knowledge) from \emph{faithfulness hallucination} (deviating from
retrieved context or instructions).
For RAG-based systems such as E-Mualim, faithfulness hallucination
is the primary concern, since the knowledge base defines the
authoritative ground truth for student-facing responses.

Hallucination is particularly damaging in educational contexts.
Kasneci et al. \cite{kasneci2023chatgpt} identified hallucination
as a central risk of LLMs in education, noting that students lack
the domain expertise to detect model errors.
Jia et al. \cite{jia2024faithfulness} found that over 23\% of
LLM-generated assignment feedback contained hallucinated content,
and Steinbach et al. \cite{steinbach2025hallucinate} demonstrated
measurable increases in student confusion under hallucinated tutoring
feedback in a pre-registered RCT.

\subsection{Retrieval-Augmented Generation}

Lewis et al. \cite{lewis2020rag} introduced RAG, combining a dense
passage retriever with a sequence-to-sequence generator and
demonstrating improvements on open-domain QA benchmarks.
Gao et al. \cite{gao2024ragsurvey} surveyed subsequent work and
identified three failure modes of \emph{naive RAG}: retrieval
imprecision (irrelevant chunks contaminating the prompt), retrieval
incompleteness (missing key passages), and unfaithful generation
(the model ignoring or contradicting retrieved context).
Shi et al. \cite{shi2023distracted} showed empirically that including
irrelevant context passages causes significant LLM performance
degradation, directly motivating chunk-level quality control.

\subsection{Advanced RAG Variants}

Several approaches improve on naive RAG by making retrieval adaptive
or self-critical.
FLARE \cite{jiang2023flare} triggers retrieval mid-generation when
token-level confidence falls below a threshold, but requires
white-box access to model probability distributions.
Self-RAG \cite{asai2024selfrag} trains the LLM itself to emit
reflection tokens for on-demand retrieval and self-critique, achieving
strong results but requiring full model retraining on specialised data.
Adaptive-RAG \cite{jeong2024adaptiverag} classifies query complexity
before retrieval to route between no-retrieval, single-step, and
multi-step strategies.
CRAG \cite{yan2024crag} takes a post-retrieval, pre-generation
approach: a fine-tuned T5-large evaluator scores retrieved documents
for relevance and triggers knowledge refinement or web-search fallback
accordingly.
E-Mualim adopts CRAG's pre-generation correction philosophy and
extends it with a post-generation faithfulness check, requiring
neither model retraining nor white-box probability access.

\subsection{RAG in Educational Settings}

RAG-based tutoring systems have grown rapidly, but almost all rely
on naive RAG \cite{sarmiento2025crag}.
Sarmiento and Laur{\'i}a \cite{sarmiento2025crag} are among the
first to compare CRAG against naive RAG in a college chatbot context,
though their evaluation is preliminary.
To our knowledge, E-Mualim represents the first educational tutoring
system combining CRAG-style chunk grading with a post-generation
hallucination checker evaluated using the RAGAS framework.

% ============================================================
%  III. SYSTEM OVERVIEW
% ============================================================
\section{E-Mualim System Overview}
\label{sec:system}

E-Mualim is a speech-enabled AI tutoring assistant that answers
student curriculum questions grounded in a curated educational
knowledge base.
The system is implemented as a LangGraph pipeline with seven nodes
across four processing layers, illustrated in Fig.~\ref{fig:pipeline}.

\begin{figure}[htbp]
  \centering
  % Uncomment and update path when figures are ready:
  % \includegraphics[width=\columnwidth]{figures/pipeline.png}
  \fbox{\rule{0pt}{4.2cm}\rule{\dimexpr\columnwidth-2\fboxsep\relax}{0pt}}
  \caption{E-Mualim LangGraph pipeline (per student message). Seven
  nodes across four layers: Input (Groq Whisper STT), Retrieve
  (intent classifier node~1, RAG retriever node~2, CRAG grader
  node~3), Generate (Llama-3.3-70b generator node~4, hallucination
  checker node~5 with max-2-retry loop), and Output (response
  expander node~6, response router node~7 with engagement-driven
  simplify loop). Dashed arrows denote corrective feedback paths.}
  \label{fig:pipeline}
\end{figure}

\subsection{Knowledge Base}

The knowledge base consists of \todo{N} educational documents
covering \todo{subject domains}, indexed in MongoDB Atlas Search
using Cohere \texttt{search\_query} embeddings.
Documents are chunked into passages of \todo{chunk size} tokens with
\todo{overlap} token overlap to preserve local context across
chunk boundaries.

\subsection{Input Layer}

Student speech is transcribed via Groq Whisper STT, producing a text
query that is forwarded to the retrieve layer.

\subsection{Retrieve Layer}

\subsubsection{Intent Classifier}
The intent classifier operates in two steps, illustrated in
Fig.~\ref{fig:intent}.
Step~1 is a deterministic keyword scan requiring no LLM call.
Messages containing terms such as ``what is'', ``define'', or
``explain'' resolve immediately to the \texttt{concept} intent;
``how do'', ``show me'', or ``example'' resolve to \texttt{howto};
error-related terms resolve to \texttt{mistake}; and similar rules
cover \texttt{practice}, \texttt{flashcard}, and
\texttt{conversational} (greetings and social messages, which bypass
RAG entirely).
The majority of student messages match a rule at this step, incurring
zero inference cost.
Step~2 fires only when no keyword matches: Groq~8b
(\texttt{llama-3.1-8b-instant}) receives a short prompt requesting a
single-word classification into one of the six categories.
The resolved intent determines which \texttt{chunk\_types} the
retriever filters on---\texttt{concept} retrieves \texttt{concept}
and \texttt{analogy} chunks, \texttt{howto} retrieves
\texttt{code\_example} and \texttt{concept} chunks,
\texttt{mistake} retrieves only \texttt{common\_mistake} chunks, and
\texttt{practice} retrieves only \texttt{practice\_problem}
chunks---ensuring the retrieved context is always relevant to the
student's actual need.

\begin{figure}[htbp]
  \centering
  % \includegraphics[width=\columnwidth]{figures/intent_classifier.png}
  \fbox{\rule{0pt}{4.2cm}\rule{\dimexpr\columnwidth-2\fboxsep\relax}{0pt}}
  \caption{Intent classifier decision tree. Step~1 applies keyword
  rules for six intents (\texttt{conversational}, \texttt{concept},
  \texttt{howto}, \texttt{mistake}, \texttt{practice},
  \texttt{flashcard}) at zero LLM cost; Step~2 invokes
  \texttt{llama-3.1-8b-instant} only when no rule matches.
  The resolved intent sets \texttt{chunk\_types} for the
  downstream RAG retriever.}
  \label{fig:intent}
\end{figure}

\subsubsection{RAG Retriever}
The RAG retriever performs \texttt{\$vectorSearch} on the Atlas
index using Cohere \texttt{search\_query} embeddings, returning the
top-$k$ chunks (\(k=\)\todo{value}) filtered by the intent-specific
\texttt{chunk\_types}.
Retrieved chunks are passed to the CRAG grader
(Section~\ref{sec:method}) before reaching the generator.

\subsection{Generate Layer}

Groq 70b (\texttt{llama-3.3-70b-versatile}) serves as the primary
generator, conditioned on the graded chunk set produced by Stage~1.
Each generated response is evaluated by the hallucination checker
(Stage~2) before it is forwarded to the output layer.

\subsection{Output Layer}

After the hallucination checker passes a response, the
\emph{response expander} (node~6) performs a citation-driven
elaboration pass.
For expandable intents (\texttt{concept}, \texttt{howto},
\texttt{mistake}, \texttt{practice}), it prompts
\texttt{llama-3.1-8b-instant} to append one enrichment paragraph
drawn from the graded source chunks---adding a concrete analogy,
code walkthrough, or common-trap warning that the 70b generator
may have omitted.
The expander is skipped for conversational and flashcard intents,
fallback responses, and retry iterations to keep latency low.

The \emph{response router} (node~7) then applies three processing
steps.
First, it extracts any fenced code block from the response and
publishes it to the frontend via a LiveKit data channel so it can
be rendered visually while the spoken response plays.
Second, it trims the response to a voice-friendly length (900
characters for expandable intents, 400 for conversational and
flashcard).
Third, it performs a single in-memory lookup on the
\texttt{SessionStateStore}:
\texttt{session\_state.get\_avg\_engagement(session\_id)}.
If the rolling average of the last five engagement readings is below
0.5, the router sets \texttt{simplify=True} in the graph state and
sends the response back to the generator with an instruction to use
simpler language and a real-world analogy
(Section~\ref{subsec:engagement}).
The final grounded response is synthesised to speech via Azure TTS
with a Bey avatar.

\subsubsection{Behavior Module (Asynchronous)}
The behavior module runs entirely independently of the LangGraph
pipeline on a 2-second interval.
Each cycle it captures a camera frame (640$\times$480), runs
DeepFace for face detection and emotion recognition, and computes an
engagement score in $[0, 1]$.
The score is written to \texttt{SessionStateStore}---a plain Python
dict in memory---via \texttt{session\_state.update()}.
The store is the \emph{only} coupling between the two subsystems.
The LangGraph pipeline never waits for the behavior loop, and the
behavior loop never waits for LangGraph; the router's dict lookup
completes in under one microsecond.

% Behavior loop diagram removed — the coupling is described in prose
% above and is secondary to the paper's CRAG contribution.

% ============================================================
%  IV. TWO-STAGE CORRECTIVE PIPELINE
% ============================================================
\section{Two-Stage Corrective Pipeline}
\label{sec:method}

Naive RAG hallucinates for two distinct reasons: the generator
receives noisy, irrelevant context (a retrieval problem), or it
drifts beyond its source material despite receiving good context
(a generation problem).
E-Mualim addresses both failure modes with two sequential corrective
stages, summarised in Fig.~\ref{fig:crag_grader}.

\begin{figure}[htbp]
  \centering
  % Uncomment and update path when figures are ready:
  % \includegraphics[width=\columnwidth]{figures/crag_grader.png}
  \fbox{\rule{0pt}{4.2cm}\rule{\dimexpr\columnwidth-2\fboxsep\relax}{0pt}}
  \caption{CRAG grader detail (Stage~1). All retrieved chunks are
  scored 0--10 in a single batched call to
  \texttt{llama-3.1-8b-instant}. Chunks scoring $\geq 6$ are
  retained; others are discarded. If all chunks are discarded,
  \texttt{needs\_fallback = true} is set and the generator falls back
  to its parametric Python knowledge.}
  \label{fig:crag_grader}
\end{figure}

\subsection{Stage 1: CRAG Chunk Grader}

\subsubsection{Motivation}
Standard RAG passes all top-$k$ chunks to the generator regardless
of their relevance to the specific query.
A vector similarity score reflects general semantic proximity but
does not guarantee that a chunk actually contains information
needed to answer the question.
We insert a dedicated relevance evaluation step between retrieval
and generation.

\subsubsection{Relevance Scoring}
All retrieved chunks are scored in a \emph{single} batched LLM call
rather than one call per chunk.
The prompt lists each chunk as a numbered excerpt and asks
\texttt{llama-3.1-8b-instant} to return a comma-separated list of
relevance scores (0--10):
\begin{quote}
\textit{``Rate the relevance of each text excerpt below to the
student question on a scale of 0--10. \ldots\ Reply with ONLY the
scores as a comma-separated list in order, e.g.: 8,3,7,5,9''}
\end{quote}
This batched, prompt-based scoring requires no fine-tuning on
labelled relevance data, in contrast to the T5-large evaluator in
the original CRAG \cite{yan2024crag}, and reduces the number of
inference calls from $k$ to one, making it practical for
resource-constrained educational deployments.

\subsubsection{Threshold Filtering}
Chunks scoring $\geq 6$ are added to the \texttt{graded\_chunks}
list and forwarded to the generator.
Chunks scoring $< 6$ are discarded.
The threshold of 6 was selected empirically on a development set
of \todo{N} query--document pairs.

\subsubsection{Fallback Behaviour}
If all $k$ retrieved chunks fall below the threshold, the flag
\texttt{needs\_fallback = true} is set.
Rather than generating from an empty context window---which
guarantees hallucination---the generator receives a modified prompt
that instructs it to answer from its parametric Python knowledge
when the question falls within the programming domain, and to
redirect the student only when the question is entirely unrelated
to programming.
This differs from the original CRAG \cite{yan2024crag}, which
triggers a web-search fallback; E-Mualim instead leverages the
base LLM's strong coverage of Python fundamentals while
constraining the scope to the tutoring domain.

\subsection{Stage 2: Post-Generation Hallucination Checker}

\subsubsection{Motivation}
Even with high-quality graded context, a generative LLM may
introduce information not present in the retained chunks.
Stage 2 provides a second, independent check on output faithfulness.

\subsubsection{Verification}
After each generator call, \texttt{llama-3.1-8b-instant} evaluates
the (graded context, generated response) pair with the prompt:
\begin{quote}
\textit{``Does the tutor answer contain information that contradicts
or goes significantly beyond the provided course material? Reply
with YES (if it goes beyond) or NO (if it stays within material).
One word only.''}
\end{quote}
A \texttt{NO} verdict (stays within material) passes the response to
the output layer.
A \texttt{YES} verdict triggers a retry, re-prompting the generator
with an explicit instruction to answer only from the provided context.
The hallucination checker is skipped for \texttt{conversational}
intents, fallback responses, and queries that have already exhausted
the retry budget.

\subsubsection{Retry Policy}
A maximum of two retries is permitted per query to bound response
latency.
If the response remains ungrounded after two retries, the most
recent response is forwarded with a confidence caveat appended.

\subsubsection{Retry Rate as a Diagnostic}
The fraction of queries triggering at least one retry serves as an
indirect measure of context quality: under CRAG-graded context,
retry rate should decrease relative to naive RAG, since cleaner
context reduces generator drift.
We report this metric alongside RAGAS scores.

\subsection{Engagement-Driven Simplification}
\label{subsec:engagement}

A third corrective mechanism specific to E-Mualim addresses response
\emph{accessibility} rather than factual grounding.
When the response router reads an \texttt{avg\_engagement} below 0.5
from the \texttt{SessionStateStore}, it sets \texttt{simplify=True}
and routes the verified response back to the generator with the
additional instruction: \emph{``Rewrite using simpler language and
include a real-world analogy.''}
This retry does not re-invoke the hallucination checker, since
simplification is constrained to rephrase already-verified content
and cannot introduce new factual claims.
Only one simplification retry is permitted per query; if engagement
remains low after simplification, the simplified response is delivered
as-is.
This loop operates only when the behavior module has written at least
one engagement reading to the store during the current session;
on the first message of a new session it is skipped.

% ============================================================
%  V. EXPERIMENTAL SETUP
% ============================================================
\section{Experimental Setup}
\label{sec:experiments}

\subsection{Test Set}

We constructed a curated test set of \todo{N} question--answer pairs
covering \todo{subject domains} drawn from E-Mualim's knowledge base.
Questions were distributed across five types to ensure coverage of
all intent-retrieval paths and CRAG grader decision boundaries:

\begin{itemize}
  \item \textbf{Concept} ($\sim$25\%) --- ``what is X'' questions
        retrieving \texttt{concept} and \texttt{analogy} chunks;
        tests baseline retrieval and generation quality.
  \item \textbf{Howto} ($\sim$25\%) --- ``how do I'' questions
        retrieving \texttt{code\_example} and \texttt{concept} chunks;
        tests code-grounded explanation quality.
  \item \textbf{Mistake} ($\sim$20\%) --- error-reporting questions
        retrieving \texttt{common\_mistake} chunks only; tests
        narrow-scope retrieval under the CRAG grader.
  \item \textbf{Complex} ($\sim$15\%) --- questions requiring
        information from multiple chunks across different sections;
        stresses context assembly and faithfulness.
  \item \textbf{Out-of-scope} ($\sim$15\%) --- questions about
        Python topics absent from the knowledge base (e.g., classes,
        file I/O); expected to trigger \texttt{needs\_fallback} and
        exercise the parametric-knowledge fallback path.
\end{itemize}
Ground-truth answers were written by the authors directly from
knowledge base content and reviewed for correctness.

\subsection{Baseline}

We compare E-Mualim against a \textbf{Naive RAG} baseline that uses
the identical retriever (Cohere embeddings, MongoDB Atlas Search),
the identical generator (\texttt{llama-3.3-70b-versatile}), and the
identical knowledge base, but with the CRAG grader, hallucination
checker, and response expander disabled.
All retrieved chunks are passed directly to the generator without
quality filtering.
This design isolates the contribution of the two-stage corrective
pipeline from all other system factors.

\subsection{Evaluation Metrics}

We evaluate both systems using the RAGAS framework \cite{es2024ragas},
which provides four complementary automated metrics:

\begin{itemize}
  \item \textbf{Faithfulness} --- proportion of claims in the
        generated answer supported by retrieved context.
        This is our primary hallucination metric; higher is better.
  \item \textbf{Answer Relevancy} --- how directly the answer
        addresses the question, penalising incomplete or off-topic
        responses.
  \item \textbf{Context Precision} --- signal-to-noise ratio of the
        retrieved (or graded) chunk set; measures filtering
        effectiveness.
  \item \textbf{Context Recall} --- whether all information necessary
        to answer the question was present in the retrieved context
        (requires ground-truth answers).
\end{itemize}

All metrics are in $[0, 1]$; higher values indicate better
performance.
We additionally report the \textbf{hallucination checker retry rate}
as an architecture-specific diagnostic.

% ============================================================
%  VI. RESULTS AND ANALYSIS
% ============================================================
\section{Results and Analysis}
\label{sec:results}

\subsection{Main Results}

Table~\ref{tab:main} reports RAGAS scores for the LLM-only baseline,
the naive RAG baseline, and the full E-Mualim system.

\begin{table}[htbp]
\centering
\caption{RAGAS Evaluation Results. All metrics in $[0,1]$, higher is
better. \todo{Replace all cells with actual numbers.}}
\label{tab:main}
\renewcommand{\arraystretch}{1.2}
\begin{tabular}{lcccc}
\toprule
\textbf{System} & \textbf{Faith.} & \textbf{Ans. Rel.}
                & \textbf{Ctx. Prec.} & \textbf{Ctx. Rec.} \\
\midrule
LLM-only              & \todo{} & \todo{} & ---     & ---     \\
Naive RAG             & \todo{} & \todo{} & \todo{} & \todo{} \\
\textbf{E-Mualim}     & \todo{} & \todo{} & \todo{} & \todo{} \\
\midrule
$\Delta$ (E-Mualim vs. Naive RAG) & \todo{} & \todo{} & \todo{} & \todo{} \\
\bottomrule
\end{tabular}
\end{table}

\todo{Write analysis paragraph: where CRAG gained most, link to
chunk grader filtering noisy context, discuss any trade-offs in
context recall from aggressive filtering.}

\subsection{Ablation Study}

Table~\ref{tab:ablation} decomposes the independent contribution of
each corrective stage.

\begin{table}[htbp]
\centering
\caption{Ablation Study. \todo{Replace all cells with actual numbers.}}
\label{tab:ablation}
\renewcommand{\arraystretch}{1.2}
\begin{tabular}{lccccc}
\toprule
\textbf{Configuration} & \textbf{Faith.} & \textbf{Ans. R.}
  & \textbf{Ctx. P.} & \textbf{Ctx. R.} & \textbf{Retry\%} \\
\midrule
Naive RAG                & \todo{} & \todo{} & \todo{} & \todo{} & \todo{} \\
+ CRAG grader only       & \todo{} & \todo{} & \todo{} & \todo{} & \todo{} \\
+ Hall.\ checker only    & \todo{} & \todo{} & \todo{} & \todo{} & \todo{} \\
\textbf{Full E-Mualim}   & \todo{} & \todo{} & \todo{} & \todo{} & \todo{} \\
\bottomrule
\end{tabular}
\end{table}

\todo{Write ablation analysis: quantify contribution of each stage
independently. Interpret retry rate reduction from naive RAG to CRAG
context as evidence that cleaner context reduces generator drift.}

\subsection{Hallucination Checker Retry Rate}

\todo{Write paragraph: report retry rate under naive RAG context vs.
CRAG-graded context. A lower retry rate under CRAG confirms that the
two stages are complementary: the grader reduces retrieval noise,
which in turn reduces post-generation correction burden.}

\subsection{Qualitative Analysis}

Table~\ref{tab:qual} presents a representative example in which
naive RAG produced a hallucinated response and E-Mualim's corrective
pipeline produced a grounded answer.

\begin{table}[htbp]
\centering
\caption{Qualitative Comparison Example.
\todo{Replace with a real query--response pair from your test set.}}
\label{tab:qual}
\renewcommand{\arraystretch}{1.3}
\begin{tabularx}{\columnwidth}{p{0.14\columnwidth}
                               X X}
\toprule
\textbf{Query} & \textbf{Naive RAG} & \textbf{E-Mualim} \\
\midrule
\todo{question here}
& \todo{hallucinated answer — include the incorrect claim}
& \todo{grounded answer — cite source passage} \\
\bottomrule
\end{tabularx}
\end{table}

% ============================================================
%  VII. DISCUSSION
% ============================================================
\section{Discussion}
\label{sec:discussion}

\subsection{Why Two Stages?}

The ablation study (Table~\ref{tab:ablation}) demonstrates that
the CRAG grader and the hallucination checker address distinct
failure modes and that their combination produces greater
hallucination reduction than either stage alone.
The grader eliminates context noise before generation, reducing the
probability that the generator has access to misleading material.
The hallucination checker provides a safety net for the residual
cases in which the generator introduces information not present in
the graded context.
Together they implement a defence-in-depth strategy appropriate for
the high-stakes requirements of educational Q\&A.

\subsection{Limitations}

Several limitations should be noted.
First, the relevance threshold ($\geq 6$) was tuned on a development
set specific to E-Mualim's knowledge base; optimal thresholds may
differ across subject domains or chunk sizes.
Second, on the critical path each query invokes
\texttt{llama-3.1-8b-instant} up to three times---once for batched
chunk grading, once for hallucination checking, and once for
response expansion---introducing additional latency relative to
naive RAG.
\todo{Report actual latency overhead in milliseconds if measured.}
Third, unlike the original CRAG \cite{yan2024crag}, E-Mualim does not
implement a web-search fallback for out-of-scope queries; the system
instead falls back to the generator's parametric Python knowledge.
While this leverages the base LLM's strong coverage of Python
fundamentals, it reintroduces the risk of parametric hallucination
for topics at the boundary of the model's training data.

\subsection{Implications for Educational AI}

Our results suggest that a two-stage corrective pipeline can be
implemented with small, freely available LLMs
(\texttt{llama-3.1-8b-instant} via Groq) without any fine-tuning,
making the architecture accessible to resource-limited institutions.
The batched scoring approach (one LLM call for all $k$ chunks)
further reduces cost relative to per-chunk evaluation.
For out-of-scope queries, E-Mualim's parametric-knowledge fallback
provides reasonable coverage of the Python domain, though it
sacrifices the grounding guarantees of the primary pipeline.

% ============================================================
%  VIII. CONCLUSION
% ============================================================
\section{Conclusion}
\label{sec:conclusion}

We presented E-Mualim, an educational tutoring system that mitigates
LLM hallucination through a two-stage corrective RAG pipeline.
The first stage filters retrieved chunks using a prompt-based
relevance grader before generation; the second verifies generated
responses against source content and retries when grounding fails.
\todo{Insert summary sentence with key result numbers here.}
Ablation results confirm that both stages contribute independently,
and the combined architecture yields the strongest hallucination
reduction.
Our work establishes a practical, reproducible design for safe
deployment of LLMs in educational Q\&A systems and provides an
empirical foundation for future work on adaptive thresholds,
multilingual support, and integration of curated web-search fallback
for open-ended student queries.

% ============================================================
%  REFERENCES
% ============================================================
\bibliographystyle{IEEEtran}
\bibliography{references_ieee}

\end{document}
